\section{Translation-invariant canonical form}

\begin{theorem}[Translation-invariant canonical form]%
    \label{thm:pgvwc07_ti_canonical_form}
    \lean{MPSTensor.exists_pgvwc07_normalized_exact_form_after_rescaling_allow_empty}
    \leanok
    % Source: PerezGarcia2007Matrix, Theorem 4.
    % Block decomposition, three block conditions, and bond-dimension bound.
    % Proof order to preserve: spectral-radius normalization and full-rank
    % fixed-point gauge; invariant support projection from a singular positive
    % fixed point; finite-ring trace split into supported/orthogonal summands;
    % iteration via a non-scalar fixed point; dual fixed point and final unitary
    % diagonalization. Derive the invariant support before any abstract
    % block-splitting lemma.
    Let $A$ be a translation-invariant MPS representation of bond
    dimension~$D$. There is a positive scalar $s$ and a finite family of
    blocks $A_k$, possibly empty, such that the globally rescaled tensor
    $s^{-1}A$ has the same MPV coefficients on every positive-length ring as
    $\bigoplus_k \mu_k A_k^i$.
    The weights are positive and satisfy $1\ge |\mu_k|$; if the family is
    nonempty, then some weight has norm~$1$. For each block,
    \begin{align}
        \sum_i A_k^i(A_k^i)^\dagger
        &= \Id,
        \notag\\
        \sum_i (A_k^i)^\dagger \Lambda_k A_k^i
        &= \Lambda_k.
        \label{eq:canonical_pgvwc_block_conditions}
    \end{align}
    Here $\Lambda_k$ is diagonal, positive, and full-rank, and each block has
    positive bond dimension. Moreover, the only fixed points of
    $\E_{A_k}(X)=\sum_i A_k^iX(A_k^i)^\dagger$ are the scalar multiples of
    $\Id$: for every $X$, the identity $\E_{A_k}(X)=X$ implies that there
    exists $c\in\C$ such that $X=c\Id$.
    The total bond dimension of the decomposition is at most~$D$.

    This is the positive-length form of
    \cite[Theorem~4]{PerezGarcia2007Matrix}, with the global
    spectral-radius normalization step made explicit.
\end{theorem}

\begin{remark}[Source proof order]\label{rem:pgvwc07_ti_canonical_source_order}
    The proof follows
    \cite[Theorem~4]{PerezGarcia2007Matrix} in order: normalize the
    spectral radius and use a full-rank positive fixed point to obtain the
    unital gauge; derive the invariant support projection from a singular
    positive fixed point by the positivity contradiction; split the finite-ring
    trace into the supported and orthogonal summands; iterate, using a
    non-scalar fixed point to split further until each block has only scalar
    fixed points; and finally repeat the fixed-point argument for the dual map
    and diagonalize the resulting full-rank positive fixed point by a unitary
    gauge. The primitive-block results below are used only after their
    hypotheses have been obtained from this argument.
\end{remark}

\begin{remark}[CPSV blocked canonical form (source statement)]%
    \label{rem:cpgsv_canonical_form_source}
    % Maintainer source anchors:
    % Papers/1606.00608/MPDO-22-12-17-2.tex: blocking removes periodic peripheral
    % components; normal-tensor and canonical-form definition eq:II_CF1 with the
    % weight normalization; any tensor becomes canonical after blocking. The
    % phase-class quotient used by the BNT supplier is prop:char-BNT, restated in
    % Appendix A. The two-layer expansion is eq:II_ABasicTensors/decBSV.
    % Papers/2011.12127/TN-Review-main.tex: invariant-subspace reduction; period
    % removal; def:4:normal-tensor-mps and eq:II_CF1; BNT
    % definition/characterization; two-layer BNT/copy display.
    The blocked canonical-form reduction of
    \cite[Section~2.3]{Cirac2016MPDO_arXiv} and
    \cite[Section~IV.A]{Cirac2021Matrix} states that an arbitrary
    translation-invariant MPS tensor becomes, after blocking by some period
    $p\ge 1$, the direct sum of an all-zero summand and a weighted family of
    normal tensors $\bigoplus_k \mu_k B_k$ with $|\mu_k|\le 1$ and at least one
    $|\mu_k|=1$. The closest proved results here are
    Theorem~\ref{thm:tp_primitive_blockdecomp} and
    Theorem~\ref{thm:paperbnt_supplier_after_blocking}; the source upper
    normalization $|\mu_k|\le 1$ with a unit-modulus witness is the one
    remaining input, recorded in
    Remark~\ref{rem:arbitrary_input_reduction_gap}.
\end{remark}

\section{Normal tensors and canonical forms}\label{sec:cpsv_canonical_defs}

This section records the CPSV definitions of \emph{normal tensor} and
\emph{canonical form} from~\cite{Cirac2016MPDO_arXiv}. The basis of normal
tensors is stated in Chapter~\ref{ch:bnt}. Write
\begin{align}
    \sigma_\partial(\E)
    &=
    \{\lambda\in\sigma(\E):|\lambda|=r(\E)\}
    \label{eq:canonical_peripheral_spectrum}
\end{align}
for the peripheral spectrum of a transfer map $\E$, where $r(\E)$ is its
spectral radius. The stronger predicates used later
(Definition~\ref{def:is_normal_canonical_form} and the canonical-form predicate
of Chapter~\ref{ch:bnt}) add left-canonical normalization and spectral
primitivity. The weight moduli are taken non-increasing, not strictly
decreasing; equal-modulus and repeated-copy sectors are retained and compared by
the coefficient argument of Chapter~\ref{ch:bnt}.

\begin{definition}[CPSV normal tensor (NT)]\label{def:nt_cpsv}
    \lean{MPSTensor.IsNormalTensor}
    \leanok
    A tensor $A$ of physical dimension~$d$ and bond dimension~$D$ is a
    \emph{normal tensor} if, after the spectral-radius normalization
    of~\cite{Cirac2016MPDO_arXiv}, (i)~$A$ admits no nontrivial invariant
    orthogonal projection and (ii)~the associated CPM
    $\E_A(X)=\sum_i A^iX(A^i)^\dagger$ has spectral radius
    $r(\E_A)=1$ and peripheral spectrum $\sigma_\partial(\E_A)=\{1\}$.
\end{definition}

\begin{remark}
    Condition~(i) cannot be dropped. A reducible tensor may have peripheral
    spectrum $\{1\}$ while still admitting a nontrivial invariant projection,
    so condition~(ii) alone does not characterize normality.
\end{remark}

The CPSV \emph{canonical form} (CF) of a tensor $A$ is the normal-block
decomposition $A^i\cong\bigoplus_{k=1}^r\mu_kA_k^i$, with weights normalized by
$|\mu_k|\le 1$ and with at least one $|\mu_k|=1$; see
\cite[Section~2.3]{Cirac2016MPDO_arXiv}. In the statements below, the structural
part is supplied by the stronger normal canonical form of
Definition~\ref{def:is_normal_canonical_form}, and the global weight
normalization is invoked as a separate source hypothesis where needed.

\begin{remark}[Relation to the strong predicates]
    A tensor $A$ which is irreducible
    (Definition~\ref{def:irreducible_tensor}), has transfer map of spectral
    radius one, and has primitive transfer map is normal in the sense of
    Definition~\ref{def:nt_cpsv}. The reverse implications from the stronger
    predicates of Definition~\ref{def:is_normal_canonical_form} and of
    Chapter~\ref{ch:bnt} require the fundamental-theorem step or the
    algebraic-to-spectral normality equivalence, established later.
\end{remark}

\subsection{Normalization conventions}\label{sec:transfer_map_normalization}

\begin{remark}[Scalar normalization conventions]
    The reductions use four scalar facts: scaling preserves injectivity; the
    transfer map under a scalar~$\zeta$ rescales by $|\zeta|^2$; unit-modulus
    phase rescaling preserves $\sum_i(A^i)^\dagger A^i=\Id$; and the MPV of a
    weighted block-diagonal tensor factorizes as
    $\mu_k^NV^{(N)}(A_k)_\sigma
    =\eta_k^N|\mu_k|^NV^{(N)}(A_k)_\sigma$, with
    $\eta_k:=\mu_k/|\mu_k|$.

    Left-canonical normalization is required before the blocks can be compared.
    Without it, the implication
    $\mc{V}(\bigoplus_k\mu_kA_k)=\mc{V}(\bigoplus_k\mu_kB_k)
    \Rightarrow\forall k,\ \mc{V}(A_k)=\mc{V}(B_k)$
    already fails for two scalar blocks: for $\mu=(1,2)$, $A=(1,3/2)$, and
    $B=(3,1/2)$, the two direct sums generate identical MPV families but the
    first blocks differ at $N=1$. Left-canonical normalization removes this
    particular ambiguity; the sectors are then matched, up to a permutation,
    by the exact fixed-length linear independence of Chapter~\ref{ch:bnt},
    which requires no ordering of the weight moduli.
\end{remark}

\section{Projector criteria and decomposition auxiliaries}

\begin{definition}[Invariant-projector closure]
    \label{def:invariant_projector_closure}
    \lean{MPSTensor.HasInvariantProjectorClosure}
    \leanok
    \uses{def:mps_tensor}
    An MPS tensor $A$ has \emph{invariant-projector closure} if every
    orthogonal projection $P$ satisfying $PA^i=PA^iP$ for all~$i$ also
    satisfies $A^iP=PA^iP$ for all~$i$.
    This is the projector hypothesis in the sufficient condition for canonical
    form of~\cite[Section~2]{Cirac2016MPDO_arXiv}.
\end{definition}

\begin{theorem}[Invariant projections are reducing under projector closure]
    \label{thm:invariant_projector_closure_commutes}
    \lean{MPSTensor.commutes_of_hasInvariantProjectorClosure_of_lowerZero}
    \leanok
    \uses{def:invariant_projector_closure}
    Suppose that $A$ has invariant-projector closure and that $P$ is an
    orthogonal projection satisfying $(\Id-P)A^iP=0$ for all~$i$. Then $P$
    reduces every tensor matrix: $PA^i=A^iP$ for all~$i$.
\end{theorem}

\begin{proof}
    \leanok
    \uses{def:invariant_projector_closure}
    The lower-corner relation gives
    $(\Id-P)A^i=(\Id-P)A^i(\Id-P)$. Applying projector closure to $\Id-P$
    gives $A^i(\Id-P)=(\Id-P)A^i(\Id-P)$, hence the upper corner
    $PA^i(\Id-P)$ also vanishes. Therefore $PA^i=PA^iP=A^iP$.
\end{proof}

\begin{theorem}[Invariant-projector closure gives a reducing projection]
    \label{thm:exists_reducing_projection_of_invariant_projector}
    \lean{MPSTensor.exists_reducing_projection_of_hasInvariantProj}
    \leanok
    \uses{def:has_inv_proj, def:invariant_projector_closure}
    If $A$ has invariant-projector closure and admits a nontrivial invariant
    projection, then it admits a nontrivial orthogonal projection commuting
    with every $A^i$.
\end{theorem}

\begin{proof}
    \leanok
    \uses{thm:invariant_projector_closure_commutes}
    Apply Theorem~\ref{thm:invariant_projector_closure_commutes} to the
    projection supplied by Definition~\ref{def:has_inv_proj}.
\end{proof}

\begin{lemma}[Support isometry of an orthogonal projection]
    \label{lem:projection_support_isometry}
    \lean{IsOrthogonalProjection.exists_support_isometry}
    \leanok
    \uses{def:orthogonal_projection_channel}
    Every orthogonal projection $P$ on $\C^D$ factors through an isometry:
    there are $n\in\N$ and $V\in\C^{D\times n}$ such that
    \begin{align}
        V^\dagger V
        &= \Id,
        \notag\\
        VV^\dagger
        &= P,
        \notag\\
        n
        &= \tr(P).
        \label{eq:canonical_projection_isometry}
    \end{align}
\end{lemma}

\begin{proof}
    \leanok
    Since $P$ is Hermitian and idempotent, its eigenvalues are $0$ and $1$.
    Writing $n=\tr(P)$, unitary diagonalization gives
    \begin{align}
        P
        &=
        U\diag(\underbrace{1,\ldots,1}_{n},0,\ldots,0)U^\dagger.
        \notag
    \end{align}
    The matrix $V\in\C^{D\times n}$ formed by the first $n$ columns of $U$
    satisfies (\ref{eq:canonical_projection_isometry}).
\end{proof}

\begin{lemma}[Corner tensors inherit projector closure]
    \label{lem:projector_closure_corner}
    \lean{MPSTensor.hasInvariantProjectorClosure_compress_of_commutes}
    \leanok
    \uses{def:invariant_projector_closure, def:orthogonal_projection_channel}
    Let $A$ have invariant-projector closure, and let $V$ be an isometry,
    $V^\dagger V=\Id$, whose range projection $VV^\dagger$ commutes with every
    $A^i$. Then the corner tensor $B^i=V^\dagger A^iV$ again has
    invariant-projector closure.
\end{lemma}

\begin{proof}
    \leanok
    \uses{def:invariant_projector_closure}
    Let $Q$ be an orthogonal projection with $QB^i=QB^iQ$ for all~$i$.
    Then $VQV^\dagger$ is an orthogonal projection. Commutation of
    $VV^\dagger$ with each $A^i$ gives $V^\dagger A^i=B^iV^\dagger$, so
    left-invariance of $Q$ yields
    \begin{align}
        (VQV^\dagger)A^i
        &=VQB^iV^\dagger
        =VQB^iQV^\dagger
        =(VQV^\dagger)A^i(VQV^\dagger).
        \notag
    \end{align}
    Projector closure for $A$ yields
    $A^i(VQV^\dagger)=(VQV^\dagger)A^i(VQV^\dagger)$, and compressing both
    sides between $V^\dagger$ and $V$ returns $B^iQ=QB^iQ$.
\end{proof}

\begin{lemma}[Matrix product vectors split along an isometric partition]
    \label{lem:mpv_partition_isometries}
    \lean{MPSTensor.sameMPV₂_toTensorFromBlocks_of_isometry_family}
    \leanok
    \uses{def:same_mpv2, def:block_diagonal_tensor}
    Let $V_1,\ldots,V_r$ be isometries, $V_k^\dagger V_k=\Id$, with
    $\sum_kV_kV_k^\dagger=\Id$, and suppose $A^iV_k=V_kB_k^i$ for some
    tensors $B_1,\ldots,B_r$. Then
    $\mc{V}(A)=\mc{V}(\bigoplus_{k=1}^rB_k)$.
\end{lemma}

\begin{proof}
    \leanok
    The intertwining relation propagates from letters to words, so
    $A^wV_k=V_kB_k^w$ for every word~$w$. Inserting the partition of unity
    $\sum_kV_kV_k^\dagger=\Id$ into the trace and cycling gives
    \begin{align}
        \tr(A^w)
        &=\sum_{k=1}^r\tr\!\left(V_k^\dagger A^wV_k\right)
        =\sum_{k=1}^r\tr(B_k^w),
        \notag
    \end{align}
    which is the matrix product vector of the direct sum.
\end{proof}

\begin{theorem}[Decomposition into irreducible corners under projector closure]
    \label{thm:projector_closure_irreducible_corners}
    \lean{MPSTensor.exists_irreducible_blockDecomp_with_isometry_of_hasInvariantProjectorClosure}
    \leanok
    \uses{def:invariant_projector_closure, def:irreducible_tensor,
        def:same_mpv2, def:block_diagonal_tensor}
    Suppose $A$ has invariant-projector closure. Then there are isometries
    $V_1,\ldots,V_r$ with $V_k^\dagger V_k=\Id$, pairwise orthogonal ranges,
    and $\sum_kV_kV_k^\dagger=\Id$, such that $A^iV_k=V_kA_k^i$ for the
    corner tensors $A_k^i=V_k^\dagger A^iV_k$, every corner tensor $A_k$ is
    irreducible, and $\mc{V}(A)=\mc{V}(\bigoplus_{k=1}^rA_k)$.
\end{theorem}

\begin{proof}
    \leanok
    \uses{def:has_inv_proj, thm:invariant_projector_closure_commutes,
        lem:projection_support_isometry, lem:projector_closure_corner,
        lem:mpv_partition_isometries}
    Proceed by strong induction on the bond dimension. If $A$ is irreducible,
    the single corner $V=\Id$ suffices. Otherwise, $A$ has a nontrivial
    invariant projection $P$, which commutes with every $A^i$ by
    Theorem~\ref{thm:invariant_projector_closure_commutes}. Factoring
    $P=V_1V_1^\dagger$ and $\Id-P=V_2V_2^\dagger$ through support isometries
    by Lemma~\ref{lem:projection_support_isometry} splits $A$ into two corner
    tensors of strictly smaller bond dimension. For $j=1,2$, commutation gives
    $A^iV_j=V_j(V_j^\dagger A^iV_j)$, and each corner inherits
    invariant-projector closure by Lemma~\ref{lem:projector_closure_corner}.
    The two recursively obtained families are combined. Ranges from different
    corners are orthogonal because
    \begin{align}
        V_1^\dagger V_2
        &=V_1^\dagger P(\Id-P)V_2
        =0.
        \notag
    \end{align}
    Equality of matrix product vectors follows from
    Lemma~\ref{lem:mpv_partition_isometries}.
\end{proof}

\begin{definition}[No nontrivial periodic vectors]
    \label{def:no_periodic_vectors}
    \lean{MPSTensor.HasNoPeriodicVectors}
    \leanok
    \uses{def:mps_tensor, def:transfer_map, def:irreducible_tensor}
    An MPS tensor $A$ has \emph{no nontrivial $p$-periodic vectors} if, for
    every isometry $V$ with $V^\dagger V=\Id$ and $A^iV=VB^i$ for all~$i$,
    with $B$ irreducible (equivalently, for every restriction $B$ of $A$ to a
    minimal invariant subspace), and for every $\rho>0$ and $r>0$ with
    $\E_B(\rho)=r\rho$, every eigenvalue of $\E_B$ of modulus~$r$ equals~$r$.
    The eigenvalue $r$ is the spectral radius of $\E_B$
    (Theorem~\ref{thm:spectral_radius_pf_irred}), so the condition states
    that no block of the invariant-subspace decomposition
    of~\cite[Section~2]{Cirac2016MPDO_arXiv}, rescaled to spectral radius
    one, has a peripheral eigenvalue $e^{2\pi iq/p}\neq1$. Eigenvectors with
    such eigenvalues are the $p$-periodic vectors
    of~\cite[Section~2]{Cirac2016MPDO_arXiv}.
\end{definition}

\begin{lemma}[Transfer map of an intertwined corner]
    \label{lem:transfer_map_corner_intertwine}
    \lean{MPSTensor.transferMap_conj_of_intertwine}
    \lean{MPSTensor.hasEigenvalue_transferMap_of_intertwine}
    \leanok
    \uses{def:transfer_map}
    Let $V$ satisfy $A^iV=VB^i$ for all~$i$. Then, for every $X$,
    \begin{align}
        \E_A(VXV^\dagger)
        &=V\E_B(X)V^\dagger.
        \label{eq:canonical_corner_transfer}
    \end{align}
    If moreover $V^\dagger V=\Id$, then every eigenvalue of $\E_B$ is an
    eigenvalue of $\E_A$.
\end{lemma}

\begin{proof}
    \leanok
    Expanding the left-hand side of (\ref{eq:canonical_corner_transfer})
    termwise, the intertwining $A^iV=VB^i$ and its adjoint
    $V^\dagger(A^i)^\dagger=(B^i)^\dagger V^\dagger$ move both factors past
    $V$ and $V^\dagger$. For the eigenvalue statement, if
    $\E_B(X)=\mu X$ with $X\neq0$, then
    $\E_A(VXV^\dagger)=\mu VXV^\dagger$, and $VXV^\dagger\neq0$ because
    compression with the isometry gives $V^\dagger(VXV^\dagger)V=X$.
\end{proof}

\begin{lemma}[Perron eigenpair of a nonzero irreducible tensor]
    \label{lem:irreducible_transfer_perron_pair}
    \lean{MPSTensor.exists_posDef_transferMap_eigenvector_of_irreducible}
    \leanok
    \uses{def:irreducible_tensor, def:transfer_map}
    Let $B$ be an irreducible MPS tensor of positive bond dimension, not all
    of whose matrices vanish. Then there are $\rho>0$ and $r>0$ with
    $\E_B(\rho)=r\rho$.
\end{lemma}

\begin{proof}
    \leanok
    \uses{thm:pf_eigenvector_existence, def:irreducible_cp}
    Irreducibility of the tensor makes the transfer map an irreducible
    completely positive map, and a nonzero matrix makes it nonzero, so it
    does not annihilate any nonzero PSD matrix. The Perron--Frobenius theorem
    for positive maps (Theorem~\ref{thm:pf_eigenvector_existence}) gives a
    nonzero PSD eigenvector with eigenvalue $r>0$, and irreducibility upgrades
    it to a positive-definite one
    (\cite[Theorem~6.3(2)]{Wolf2012Quantum}).
\end{proof}

\begin{lemma}[Spectral normalization of an irreducible block]
    \label{lem:irreducible_block_spectral_normalization}
    \lean{MPSTensor.isNormalTensor_invSqrt_smul_of_unique_peripheral}
    \leanok
    \uses{def:irreducible_tensor, def:transfer_map, def:nt_cpsv}
    Let $B$ be an irreducible MPS tensor of positive bond dimension with
    $\E_B(\rho)=r\rho$ for some $\rho>0$ and $r>0$, and suppose every
    eigenvalue of $\E_B$ of modulus~$r$ equals~$r$. Then $r^{-1/2}B$ is a
    normal tensor.
\end{lemma}

\begin{proof}
    \leanok
    \uses{def:nt_cpsv, thm:spectral_radius_pf_irred}
    The rescaling identity $\E_{r^{-1/2}B}(X)=r^{-1}\E_B(X)$ gives
    $\E_{r^{-1/2}B}(\rho)=r^{-1}r\rho=\rho$, so $\rho$ is a
    positive-definite fixed point of $\E_{r^{-1/2}B}$. The rescaled map
    remains irreducible and completely positive, so
    Theorem~\ref{thm:spectral_radius_pf_irred} identifies its spectral radius
    with the eigenvalue one. The eigenvalues of modulus one of
    $\E_{r^{-1/2}B}$ are the eigenvalues of $\E_B$ of modulus~$r$, divided by
    $r$; by hypothesis the only such eigenvalue is one, so the peripheral
    eigenvalue set is $\{1\}$. Irreducibility is unchanged by a nonzero
    rescaling.
\end{proof}

\begin{lemma}[Canonical inclusion of a dependent direct summand]
    \label{lem:matrix_sigma_block_inclusion}
    \lean{Matrix.sigmaBlockInclusion}
    \lean{Matrix.sigmaBlockInclusion_isometry}
    \lean{Matrix.blockDiagonal'_mul_sigmaBlockInclusion}
    \lean{Matrix.sigmaBlockInclusion_conjTranspose_mul_blockDiagonal'}
    \lean{Matrix.sigmaBlockInclusion_compression}
    \lean{Matrix.blockDiagonal'_eq_sum_sigmaBlockInclusion}
    \lean{Matrix.blockDiagonal'_commute_sigmaBlockInclusion_rangeProjection}
    \leanok
    Let $E_k:H_k\to\bigoplus_lH_l$ be the canonical inclusion and let
    $B=\bigoplus_lB_l$. Then
    \begin{align}
        E_k^\dagger E_k
        &=\Id,
        \notag\\
        BE_k
        &=E_kB_k,
        \notag\\
        E_k^\dagger B
        &=B_kE_k^\dagger,
        \notag\\
        E_k^\dagger BE_k
        &=B_k,
        \notag\\
        B
        &=\sum_kE_kB_kE_k^\dagger.
        \label{eq:canonical_direct_summand}
    \end{align}
    In particular, $B$ commutes with the range projection
    $E_kE_k^\dagger$.
\end{lemma}

\begin{proof}\leanok
    The identities in (\ref{eq:canonical_direct_summand}) follow directly
    from the defining block-diagonal entries. The relations
    $BE_k=E_kB_k$ and $E_k^\dagger B=B_kE_k^\dagger$ imply that $B$
    commutes with $E_kE_k^\dagger$.
\end{proof}

\begin{lemma}[Canonical hypotheses descend to dependent direct summands]
    \label{lem:canonical_hypotheses_dependent_direct_summand}
    \lean{MPSTensor.projectorClosure_and_noPeriodicVectors_block_of_reindex_eq_blockDiagonal}
    \leanok
    \uses{lem:matrix_sigma_block_inclusion,
        def:invariant_projector_closure, def:no_periodic_vectors}
    Suppose that, after a change of coordinates, $A$ is the dependent direct
    sum of tensors $B_k$. If $A$ has invariant-projector closure and no
    nontrivial periodic vectors, then every $B_k$ has both properties. The
    assertion also applies to summands of dimension zero.
\end{lemma}

\begin{proof}\leanok
    Transport the canonical inclusion of the $k$-th summand back to the bond
    coordinates of $A$. This is an isometry which intertwines $B_k$ with
    $A$ on both sides, and its range projection commutes with every letter of
    $A$. Projector closure therefore descends by exact compression. A
    periodic vector of $B_k$ would embed through the same isometry as a
    periodic vector of $A$.
\end{proof}

\begin{lemma}[Closed-chain equality from an exact isometric decomposition]
    \label{lem:same_mpv_exact_isometric_decomposition}
    \lean{MPSTensor.sameMPV₂Pos_toTensorFromBlocks_of_reconstruction}
    \leanok
    \uses{def:same_mpv2_pos, def:block_diagonal_tensor}
    Suppose that isometries $V_k$ satisfy
    \begin{align}
        V_k^\dagger A^i
        &=(\mu_kB_k^i)V_k^\dagger,
        \notag\\
        A^i
        &=\sum_kV_k(\mu_kB_k^i)V_k^\dagger.
        \label{eq:canonical_exact_reconstruction}
    \end{align}
    Then $A$ and $\bigoplus_k\mu_kB_k$ have the same matrix product vectors
    at every positive length.
\end{lemma}

\begin{proof}\leanok
    Expand the first letter by the reconstruction formula in
    (\ref{eq:canonical_exact_reconstruction}). In each summand, move the
    remaining word through $V_k^\dagger$ using the intertwining relation.
    Cyclicity of the trace places $V_k^\dagger V_k=\Id$ next to the weighted
    block word, leaving precisely its closed-chain trace.
\end{proof}

\begin{theorem}[Sufficient condition for canonical form]
    \label{thm:projector_closure_canonical_form}
    \lean{MPSTensor.%
        exists_normalTensor_blockDecomp_with_isometry_of_hasInvariantProjectorClosure}
    \lean{MPSTensor.exists_normalTensor_blockDecomp_sameMPV₂Pos_of_hasInvariantProjectorClosure}
    \leanok
    \uses{def:invariant_projector_closure, def:no_periodic_vectors,
        def:nt_cpsv, def:same_mpv2_pos, def:block_diagonal_tensor}
    Suppose $A$ has invariant-projector closure and no nontrivial
    $p$-periodic vectors. Then there are normal tensors
    $A_1,\ldots,A_r$ of positive bond dimensions $D_1,\ldots,D_r$, with
    $\sum_kD_k\le D$, and nonzero weights $\mu_1,\ldots,\mu_r$ such that
    \begin{align}
        \mc{V}^+(A)
        &=
        \mc{V}^+\!\left(\bigoplus_{k=1}^r\mu_kA_k\right).
        \label{eq:canonical_projector_decomp}
    \end{align}
    Moreover, there are isometries $V_k:\C^{D_k}\to\C^D$ with pairwise
    orthogonal ranges such that
    \begin{align}
        A^iV_k
        &=V_k(\mu_kA_k^i),
        \notag\\
        A^i
        &=\sum_kV_k(\mu_kA_k^i)V_k^\dagger.
        \label{eq:canonical_projector_isometries}
    \end{align}
    This is the sufficient condition for canonical form
    of~\cite[Section~2]{Cirac2016MPDO_arXiv}, with that paper's convention
    $\sum_kD_k\le D$ allowing zero blocks: a corner of $A$ carrying the
    zero tensor cannot be rescaled to spectral radius one and is omitted
    because it vanishes at every positive length.
\end{theorem}

\begin{proof}
    \leanok
    \uses{thm:projector_closure_irreducible_corners,
        lem:irreducible_transfer_perron_pair,
        lem:irreducible_block_spectral_normalization}
    Decompose $A$ into irreducible corners $B_k$ along isometries $V_k$
    by Theorem~\ref{thm:projector_closure_irreducible_corners}; every matrix
    product vector coefficient of $A$ is the sum of those of the corners. A
    corner whose matrices all vanish contributes nothing at positive length
    and is dropped. Every remaining corner has a Perron eigenpair
    $\E_{B_k}(\rho_k)=r_k\rho_k$ by
    Lemma~\ref{lem:irreducible_transfer_perron_pair}. The absence of
    $p$-periodic vectors, applied to the corner $B_k$, gives that every
    eigenvalue of $\E_{B_k}$ of modulus~$r_k$ equals~$r_k$, so
    $A_k:=r_k^{-1/2}B_k$ is normal by
    Lemma~\ref{lem:irreducible_block_spectral_normalization}. Setting
    $\mu_k:=r_k^{1/2}$ gives $\mu_kA_k=B_k$ at the level of matrices, so for
    every word $w$ of positive length~$N$ the weight cancels:
    \begin{align}
        \mu_k^N\tr(A_k^w)
        &=
        (r_k^{1/2})^N(r_k^{-1/2})^N\tr(B_k^w)
        =
        \tr(B_k^w).
        \notag
    \end{align}
    Summing over $k$ gives (\ref{eq:canonical_projector_decomp}), while the
    corner reconstruction gives (\ref{eq:canonical_projector_isometries}).
    The retained dimensions satisfy
    $\sum_kD_k\le\sum_k\tr(V_kV_k^\dagger)=\tr(\Id)=D$.
\end{proof}

\begin{theorem}[Non-scalar fixed point gives a further split]%
    \label{thm:unital_nonscalar_fixed_point_split}
    \lean{MPSTensor.exists_singular_posSemidef_fixedPoint_of_unital_nonScalar_fixedPoint}
    \lean{MPSTensor.exists_twoBlock_decomp_of_unital_nonScalar_fixedPoint}
    \leanok
    \uses{def:transfer_map, def:same_mpv2, thm:two_block_decomp}
    Suppose $A$ is unital, $\E_A(\Id)=\Id$. If $X$ is a Hermitian fixed point
    of $\E_A$ which is not a scalar multiple of $\Id$, then there is a nonzero
    PSD fixed point $\rho$ which is not positive definite, and consequently
    $A$ admits a two-block MPV-equivalent decomposition with both block
    dimensions strictly smaller than~$D$. This is the fixed-point shift and
    further decomposition step in
    \cite[Theorem~4]{PerezGarcia2007Matrix}.
\end{theorem}

\begin{proof}\leanok
    Let $\lambda_{\max}$ be the largest eigenvalue of $X$. Then
    $\rho:=\lambda_{\max}\Id-X$ is positive semidefinite, singular, and
    nonzero because $X$ is not scalar. Unitality and the fixed-point equation
    for $X$ give $\E_A(\rho)=\rho$. Apply
    Theorem~\ref{thm:two_block_decomp}. The non-scalar hypothesis is needed
    because scalar multiples of $\Id$ are fixed by every unital map and yield
    no nontrivial support projection.
\end{proof}
